\documentclass[11pt]{article}
\usepackage[margin=1in]{geometry}
\usepackage{booktabs}
\usepackage{longtable}
\usepackage{hyperref}
\usepackage{enumitem}
\title{Independent Replication Report\\
\large ``Analysis of the complete genome sequence of \textit{Nocardia seriolae} UTF1,
the causative agent of fish nocardiosis''\\
Yasuike et al., \textit{PLoS ONE} \textbf{12}(3):e0173198 (2017)}
\author{Ollie (OpenClaw AI) --- BVBRC-100 Replication Wave, target BVBRC-63}
\date{Backfilled to 8-artifact standard: 2026-07-05}
\begin{document}
\maketitle

\section{Paper summary}
Yasuike et al.\ (2017) report the first complete reference genome for \textit{N.\ seriolae},
causative agent of yellowtail/amberjack/kingfish nocardiosis in Japanese aquaculture, and
perform comparative genomics against the four then-available complete \textit{Nocardia}
genomes (\textit{N.\ farcinica} IFM 10152, \textit{N.\ brasiliensis} HUJEG-1,
\textit{N.\ cyriacigeorgica} GUH-2, \textit{N.\ nova} SH22a). Key findings: single 8.12 Mb
circular chromosome, 68.1\% G+C, 7{,}697 predicted CDS, 2{,}745 core orthologs, 1{,}982
UTF1-unique genes enriched for mobile elements / hypothetical proteins / ABC transporters,
plus mycobacterial-style virulence-factor orthologs.

\section{Claims tested}
\begin{longtable}{clcc}
\toprule
\# & Claim & Type & Tested \\
\midrule
C1 & Circular chromosome 8{,}121{,}733 bp & Quantitative & Yes \\
C2 & G+C 68.1\% & Quantitative & Yes \\
C3 & 7{,}697 predicted proteins & Quantitative & Yes \\
C4 & 4 rRNA operons (16S/23S/5S) & Quantitative & Yes \\
C5 & 2{,}745 core orthologs (all 5 strains) & Quantitative & Yes \\
C6 & 1{,}982 UTF1-unique genes & Quantitative & Yes \\
C7 & Mobile-element enrichment vs comparators & Quant & Yes \\
C8 & Unknown-function-gene enrichment & Quant & Partial \\
C9 & ABC-transporter enrichment in unique genes & Quant & Partial \\
C10 & Mycobacterial/human-\textit{Nocardia} virulence orthologs & Qualitative & Yes \\
\bottomrule
\end{longtable}

\section{Method}
Assembly RefSeq GCF\_002356035.1 (GenBank AP017900.1; PacBio RS, 133$\times$) from NCBI.
Basic stats via Biopython. Independent re-annotation via Prokka 1.12 (Prodigal 2.6.3) on
uicgpu (8$\times$A100). Comparator RefSeq assemblies for the four named strains retrieved.
Orthology by reciprocal-best-hit BLASTP (e-value 1e-5, pid$\geq$25\%, cov$\geq$40\%): UTF1
protein is ``core'' if RBH in all 4 comparators, ``unique'' if RBH in none. Functional
categories by keyword counting across the 5 RefSeq \texttt{.faa} description lines (coarse
proxy for direction-of-enrichment). Verdict by LLM judge (Argo \texttt{argo:claude-opus-4.7})
--- no regex scoring of the verdict.

\section{Results vs paper}
\begin{longtable}{llll}
\toprule
\# & Paper & This work & Agreement \\
\midrule
C1 & 8{,}121{,}733 bp & 8{,}121{,}733 bp & EXACT \\
C2 & 68.1\% & 68.14\% & EXACT (rounding) \\
C3 & 7{,}697 & RefSeq 7{,}650 / Prokka 7{,}648 & 99.4\% \\
C4 & 4 operons (12 rRNA) & 12 rRNA / 4 each & EXACT \\
C5 & 2{,}745 core & 2{,}718 RBH & 99.0\% \\
C6 & 1{,}982 unique & 1{,}967 & 99.2\% \\
C7 & ``greater'' mobile & 127 vs 43/24/24/20 & STRONG (3.0--6.4$\times$) \\
C8 & unknown-fn enrichment & 22.3\% (2nd of 5) & PARTIAL (not distinctly higher on \%) \\
C9 & ABC in unique genes & 236 (3.3\%), sim. \% & PARTIAL (elevated absolute, similar \%) \\
C10 & virulence orthologs & Mce 21, cat/SOD 4, sidero 3, efflux 11, $\beta$-lac 1 & Present \\
\bottomrule
\end{longtable}

\section{Verdict}
\textbf{REPLICATED.} All quantitative genome-level claims (C1--C6) reproduce within $\leq$1\%
on the deposited assembly + independent re-annotation. The central comparative-genomics result
(2{,}745 core / 1{,}982 unique) reproduces to 99.0\% / 99.2\% via fresh RBH BLASTP against the
same four comparators. Mobile-element enrichment (C7) is dramatically confirmed (3--6$\times$).
Virulence orthologs (C10) present in expected classes. C8/C9 only partially supported by the
coarse keyword proxy, neither contradicted. Coverage 10/10, 8 confirmed, 2 partial, 0
contradicted. LLM judge concurs (REPLICATED, $\sim$85\% strength).

\section{Genuine critique (evidence strength, shortcuts, unverified claims)}
\begin{itemize}[leftmargin=1.4em]
\item \textbf{C8/C9 tests are structurally weaker than the paper's claims.} The paper claims
enrichment WITHIN the UTF1-unique gene set; our keyword proxy counted categories across the
WHOLE proteome, not restricted to unique genes. So ``partial'' here means ``we did not
actually test the paper's specific claim'' rather than ``we tested and it half-held.'' This is
an honest scope gap, not a confirmation.
\item \textbf{Functional categorization is regex-on-description-lines}, not COG/eggNOG/Pfam.
This is a coarse direction-of-enrichment proxy only. Real functional enrichment testing
(hypergeometric over COG categories) was not performed and could shift C7--C9 conclusions.
\item \textbf{RBH orthology thresholds (pid$\geq$25\%, cov$\geq$40\%) are loose genus-level
choices.} The 99\% agreement on core/unique counts is impressive but conditioned on these
cutoffs; a sensitivity sweep was not run, so we cannot claim the 2{,}718/1{,}967 numbers are
threshold-robust.
\item \textbf{CDS count discrepancy across sources} (paper 7{,}697; RefSeq 7{,}650; Prokka
7{,}648; and a per-species table row shows UTF1 7{,}130 in one annotation) reflects
annotation-pipeline divergence. We report all but did not reconcile which gene models differ
--- the 99.4\% agreement masks which specific ORFs are gained/lost.
\item \textbf{Virulence-factor claim (C10) is presence-only.} We confirmed orthologs of Mce,
catalase/SOD, siderophore, efflux, $\beta$-lactamase exist, but did NOT verify functional
virulence or that these are the SAME orthologs the paper highlights --- just that the classes
are represented.
\item \textbf{Single LLM judge} (Argo opus-4.7), below the 3-judge policy ideal. Defensible
given the exact quantitative agreement, but noted.
\end{itemize}

\section{Open questions}
See \texttt{report/open\_questions.json} --- five open problems: (Q1) COG/eggNOG-based
enrichment to properly test C8/C9 within the unique gene set; (Q2) RBH threshold sensitivity;
(Q3) UTF1-unique mobile-element characterization + horizontal-transfer origins; (Q4)
functional validation of the mycobacterial-style virulence orthologs in a fish-infection
model; (Q5) pan-genome expansion with the many \textit{N.\ seriolae} genomes sequenced since 2017.

\end{document}
